\DeclareMathOperator{\Frac}{Frac}
\DeclareMathOperator{\Tot}{Tot}
\title{A Complete Local Domain without\\a Small Cohen--Macaulay Module}
\author{OpenAI}
\date{September 23, 2026}
\begin{document}
\maketitle
\begin{abstract}
We construct a three-dimensional complete Noetherian normal local domain
containing $\C$, with residue field $\C$, that has no nonzero finitely generated
maximal Cohen--Macaulay module. This disproves the domain form of the small
Cohen--Macaulay module conjecture.
\end{abstract}
\tableofcontents
\clearpage

\section{Introduction}
A \emph{small Cohen--Macaulay module} over a Noetherian local ring $(R,\mathfrak m)$
is a nonzero finitely generated $R$-module $M$ such that
$\depth_R M=\dim R$. Such a module is also called a maximal Cohen--Macaulay
module. The domain form of the small Cohen--Macaulay module conjecture asks
whether every complete Noetherian local domain has one.
We prove the following negative answer.

\begin{theorem}\label{thm:main}
There exists a three-dimensional complete Noetherian normal local domain
$R$ containing $\C$, with residue field $\C$, such that no nonzero finitely
generated $R$-module has depth three.
\end{theorem}

The distinction between modules and algebras matters here. A module-finite extension domain $R\subseteq B$ that is
Cohen--Macaulay supplies such a module, but the converse need not hold. Our obstruction applies to every finite module;
it requires neither a grading nor an algebra structure.
Hochster formulated the finite-module existence problem in his early
work on Cohen--Macaulay modules and homological conjectures
\cite{Hochster1973,Hochster1975}. In dimensions at most two, a complete
local domain has a small Cohen--Macaulay module: its finite normalization
is Cohen--Macaulay \cite[Section~2]{Hochster2017}. There are also important
positive results in the graded setting. The theorem of Hartshorne,
Peskine--Szpiro, and Hochster gives graded small Cohen--Macaulay modules
for three-dimensional finitely generated nonnegatively graded domains
